\chapter{Estado da Arte}
\label{chap:estado_arte}

\section{Introdução}
\label{sec:ea_introducao}
A expansão de dispositivos da \textit{Internet of Things} (IoT) trouxe consigo desafios de segurança sem precedentes. Este capítulo apresenta uma revisão da literatura sobre as vulnerabilidades inerentes às redes IoT, a evolução da cibersegurança suportada por \textit{Machine Learning} (ML) e os constrangimentos arquiteturais de implementar estas soluções na periferia da rede (\textit{Edge Computing}).

\section{Segurança em Redes IoT e a Evolução da Deteção de Intrusões}
\label{sec:ea_seguranca_ids}
Os dispositivos IoT, como sensores e atuadores, são frequentemente desenhados com restrições ao nível do processamento, memória e consumo energético \cite{scavenging}. Esta limitação computacional inviabiliza a instalação de protocolos de segurança complexos ou soluções de antivírus tradicionais nos próprios equipamentos. Consequentemente, a superfície de ataque nestas redes é vasta, tornando-as alvos preferenciais para o recrutamento em \textit{Botnets} e para a orquestração de ataques de Negação de Serviço Distribuída (DDoS) \cite{Baccelli2019}. Adicionalmente, a heterogeneidade dos protocolos de comunicação dificulta a padronização das defesas, exigindo mecanismos que operem diretamente na camada de rede \cite{IEEE_Anomalias_2021}.

Historicamente, os Sistemas de Deteção de Intrusão (IDS) basearam-se em catálogos de assinaturas (\textit{Signature-based IDS}), comparando o tráfego com padrões conhecidos de \textit{malware}. Embora eficazes contra ameaças documentadas, estes sistemas revelam-se ineficazes perante ataques \textit{Zero-Day} ou mutações de código malicioso \cite{AlGaradi2020}. Para colmatar esta lacuna, a comunidade científica tem transitado para abordagens baseadas em anomalias suportadas por \textit{Machine Learning}. Modelos como \textit{Random Forest}, \textit{Support Vector Machines} (SVM) e Redes Neuronais Artificiais (ANN) demonstraram uma elevada capacidade para modelar o comportamento "normal" da infraestrutura e classificar desvios estatísticos como potenciais intrusões \cite{Berman2019}.

\section{Desafios Arquiteturais: Análise de Fluxos e o Paradoxo da Latência}
\label{sec:ea_desafios_arquiteturais}
A eficácia de um modelo de Inteligência Artificial depende diretamente do método de extração de características (\textit{features}) do tráfego. A Inspeção Profunda de Pacotes (\textit{Deep Packet Inspection} - DPI) analisa o conteúdo integral da mensagem (\textit{payload}), garantindo alta precisão. Contudo, o DPI levanta sérias questões de privacidade e exige um esforço computacional incompatível com \textit{routers} de baixo custo \cite{IEEE_Anomalias_2021}. Em alternativa, a análise baseada em fluxos (\textit{Flow-based Analysis}) extrai apenas metadados estatísticos da comunicação, tais como o tempo entre pacotes, a duração do fluxo e a contagem de \textit{flags} TCP. Ferramentas como o \textit{CICFlowMeter}, utilizadas para a geração de \textit{datasets} como o CICIoT2023 \cite{CICIoT2023}, provam que esta abordagem garante a preservação da privacidade e reduz drasticamente a sobrecarga de processamento.

Esta redução de carga computacional é um fator crítico quando se introduz o paradigma de \textit{Edge Computing}, que propõe a deslocação do processamento de dados da nuvem (\textit{Cloud}) para a periferia da rede \cite{Chen2019}. No entanto, a literatura atual evidencia um paradoxo na aplicação destas tecnologias: a esmagadora maioria dos estudos foca-se exclusivamente na maximização da exatidão (\textit{Accuracy}) preditiva, recorrendo a \textit{Ensembles} de árvores de decisão altamente complexos. O problema fundamental desta abordagem reside no negligenciar do "orçamento de latência". Em cenários de ataques volumétricos, o tempo que um \textit{Edge Router} demora a processar a inferência matemática de uma árvore profunda excede o tempo de chegada do pacote seguinte, resultando no bloqueio do \textit{buffer} e na negação de serviço provocada pelo próprio sistema de defesa.

\section{Posicionamento do Projeto e Comparação com a Literatura}
\label{sec:ea_posicionamento}

Em contraste com grande parte da literatura existente, o presente trabalho não procura unicamente a precisão estatística isolada, mas sim a resolução do problema da latência em ambientes de \textit{Edge Computing}. 

Para estabelecer um termo de comparação direto e fundamentar as decisões arquiteturais tomadas, é pertinente analisar trabalhos recentes com metodologias de simulação semelhantes. Destaca-se o estudo focado em \textit{SDN-Enabled IoT Anomaly Detection Using Ensemble Learning} \cite{EnsembleSDN2020}, que também recorre ao emulador Mininet para validar mecanismos de segurança numa rede \ac{IoT}. Nesse estudo, os autores propõem a utilização de \textit{Ensembles} de aprendizagem profunda (combinando múltiplas redes e árvores) colocados num controlador central para maximizar a exatidão e o \textit{F1-Score}. Embora a abordagem alcance uma precisão teórica notável na deteção de ataques dinâmicos, os autores focam-se quase exclusivamente no poder preditivo abstrato, não acautelando a sobrecarga computacional de processar modelos tão complexos em tempo real.

É exatamente nesta lacuna que este projeto se demarca. A Tabela \ref{tab:comparacao_literatura} sintetiza os compromissos assumidos neste trabalho, em oposição à abordagem tradicional focada na maximização do acerto.

Este trabalho assume que, em cenários de ataques volumétricos (como injeções massivas de pacotes UDP/TCP), utilizar \textit{Ensembles} pesados na periferia da rede resulta numa latência insustentável. Ao avaliar empiricamente o custo-benefício, o presente projeto coroa modelos mais leves, como o \ac{MLP}, sacrificando uma margem negligenciável de exatidão pura em prol da sobrevivência da rede. 

Adicionalmente, a inovação prática deste projeto reside na hibridização do sistema: ao acoplar o modelo \ac{MLP} a uma arquitetura de \textit{Blacklists} dinâmicas operada diretamente no Servidor \textit{Edge} (atuando como um nó de segurança dedicado), garante-se uma latência de processamento contínua mínima. Esta separação física e lógica de funções liberta o \textit{router} principal da pesada carga de filtragem, impedindo o bloqueio do seu \textit{buffer} e viabilizando a proteção em tempo real para infraestruturas com recursos limitados.


\begin{table}[htbp]
    \centering
    \caption{Comparação de \textit{Trade-offs}: Literatura \textit{vs.} Arquitetura Proposta.}
    \label{tab:comparacao_literatura}
    \renewcommand{\arraystretch}{1.4} % Aumenta ligeiramente o espaçamento das linhas para ficar mais elegante
    \begin{tabular}{|p{3.5cm}|p{5cm}|p{5cm}|}
        \hline
        \textbf{Critério de Avaliação} & \textbf{Artigo de Referência \cite{EnsembleSDN2020}} & \textbf{O Nosso Projeto (MLP + \textit{Blacklist})} \\
        \hline
        \textbf{Topologia de Rede} & Controlador Central SDN & Nó de \textit{Edge Computing} autónomo \\
        \hline
        \textbf{Algoritmo de IA} & \textit{Ensemble Learning} (Múltiplos modelos pesados) & \textit{Multi-Layer Perceptron} (Cálculo matricial rápido) \\
        \hline
        \textbf{Foco Principal} & Maximização absoluta da Exatidão (\textit{Accuracy}) & Equilíbrio ótimo entre Exatidão preditiva e Latência \\
        \hline
        \textbf{Mecanismo de Mitigação} & O controlador SDN processa e bloqueia de forma contínua & Arquitetura Híbrida (IA analisa o 1º pacote; \textit{Blacklist} bloqueia os restantes) \\
        \hline
        \textbf{Latência Operacional} & Alta (inviável para processamento de pacotes individuais) & Ultrabaixa ($\sim$1 ms por inferência preditiva) \\
        \hline
        \textbf{Risco em DDoS} & Risco severo de exaustão de recursos e estrangulamento da rede & Mitigação quase instintiva, poupando a capacidade do \textit{Buffer} \\
        \hline
    \end{tabular}
\end{table}
